\par\medskip
\Needspace{85mm}
\begin{center}
\tagstructbegin{tag=Figure,alttext={Best-of-N 教学成功率图。横轴为一到十六份候选，纵轴为交付正确答案的概率。独立候选且完美核验时，单次成功率四分之一，曲线经过 N 为一、二、四、八、十六时的百分之二十五、四十三点七五、六十八点三六、八十九点九九、九十九点零零。候选完全复制时，成功率始终为百分之二十五。}}
\begin{tikzpicture}
\begin{axis}[width=.88\linewidth,height=61mm,
 xmin=1,xmax=16,ymin=0,ymax=1.06,
 xtick={1,2,4,8,16},ytick={0,.25,.5,.75,1},
 xlabel={候选数 $N$},ylabel={成功率},
 tick label style={font=\footnotesize},label style={font=\small},
 axis lines=left,axis line style={draw=muted},
 grid=major,grid style={draw=linegray},
 legend style={font=\footnotesize,draw=none,fill=white,at={(.97,.18)},anchor=south east}]
\addplot[blue,very thick,domain=1:16,samples=100]{1-.75^x};
\addlegendentry{独立采样＋完美核验}
\addplot[only marks,mark=*,blue] coordinates {(1,.25) (2,.4375) (4,.68359375) (8,.89988708496) (16,.98997740424)};
\addlegendentry{手算的五个点}
\addplot[orange,thick,dashed,domain=1:16]{.25};
\addlegendentry{完全复制同一次答案}
\end{axis}
\end{tikzpicture}
\captionof{figure}{多尝试的收益取决于候选的差异和核验能力；曲线是教学假设下的计算结果}
\label{fig:best-of-n-success}
\tagstructend
\end{center}
